\subsection{کالبدشکافی ماژول \lr{phantom}}

\subsubsection*{۱. فلسفه طراحی، پارادایم معماری و تمثیل ملموس}
معماری نوآورانه \lr{phantom} با معرفی توکن‌های مجازی نهفته («توکن‌های شبح»)، تعاملات میان بخش‌های دوردست متن را بدون مادی‌سازی اتصالات جفتی مربعی برقرار می‌کند. توکن‌های شبح به عنوان پل‌های نامرئی اطلاعاتی عمل کرده و وضعیت کلی توالی را با کمترین سربار پردازشی انتقال می‌دهند.

\textbf{تمثیل ملموس و مثال کاربردی:}
ارسال سیگنال‌های راهنما \lr{(Pilot Signals)} در شبکه‌های فیبر نوری پرسرعت: برای هدایت بسته‌های اطلاعاتی و تنظیم همگام‌سازی امواج بدون اشغال طول موج‌های داده‌ای اصلی، پالس‌های مجازی شبح‌گون با انرژی بسیار کم به موازات جریان داده حرکت می‌کنند تا وضعیت خط را گزارش دهند بدون آنکه پهنای باند اصلی مصرف شود.

\subsubsection*{۲. فرمول‌بندی و استنتاج عمیق ریاضی \lr{(Mathematical Derivation)}}
دنباله اصلی $\mathbf{X} \in \mathbb{R}^{B \times N \times D}$ با $P \ll N$ توکن شبح آموزش‌پذیر $\mathbf{Z} \in \mathbb{R}^{P \times D}$ ترکیب می‌شود:
\begin{equation}
\tilde{\mathbf{X}} = [\mathbf{X} \, ; \, \mathbf{Z}] \in \mathbb{R}^{B \times (N+P) \times D}
\end{equation}
ماتریس اتصال از دو گام فرافکنی دوطرفه تشکیل می‌شود: نگاشت توکن‌های ورودی به توکن‌های شبح $\mathbf{A}_{X \to Z}$ و سپس بازتوزیع اطلاعات از توکن‌های شبح به تمام توکن‌ها $\mathbf{A}_{Z \to X}$:
\begin{equation}
\operatorname{PhantomAttn}(\mathbf{X}) = \operatorname{Softmax}\left(\frac{\mathbf{Q}_X \mathbf{K}_Z^\top}{\sqrt{d}}\right) \cdot \operatorname{Softmax}\left(\frac{\mathbf{Q}_Z \mathbf{K}_X^\top}{\sqrt{d}}\right) \mathbf{V}_X
\end{equation}
این ساختار دوگانه، پیچیدگی را به سطح کاملاً خطی $\mathcal{O}(N P d)$ کاهش می‌دهد.

\subsubsection*{۳. کالبدشکافی الگوریتمی و جریان داده \lr{(Algorithmic Tensor Trace)}}
در پیاده‌سازی سورس پایتون در فایل ماژول \lr{phantom.py}، پردازش در متد \lr{forward} به صورت مرحله‌به‌مرحله و گام‌به‌گام صورت می‌پذیرد:
\begin{enumerate}
\item الحاق تانسور توکن‌های شبح پارامتری $\mathbf{Z}$ به انتهای تانسور دنباله ورودی.
\item محاسبه پروجکشن‌های خطی برای توکن‌های اصلی و توکن‌های شبح.
\item \textbf{اجرای مرحله اول توجه فشرده:} نگاشت اطلاعات کل دنباله به توکن‌های شبح با اندازه $[B, H, P, d]$.
\item \textbf{اجرای مرحله دوم توجه:} بازپخش اطلاعات غنی‌شده از توکن‌های شبح به مواضع ورودی اصلی.
\item جداسازی تانسور توکن‌های اصلی و ارسال به پروجکشن خروجی با جمع مانده.
\end{enumerate}

\subsubsection*{۴. تحلیل همجوشی سخت‌افزاری و ترافیک حافظه \lr{(Hardware \& Roofline Analysis)}}
از دیدگاه مدل سقف کارایی \lr{(Roofline Model)} و معماری پردازنده‌های گرافیکی مدرن (نظیر کارت گرافیک \lr{NVIDIA RTX 4070}):
\begin{itemize}
\item \textbf{سازگاری کامل با کش رجیستری:} تعداد اندک توکن‌های شبح ($P=16$) اجازه می‌دهد کل بافر میانی درون ثبات‌های سریع \lr{SRAM} قرار گیرد.
\item \textbf{حذف کامل وابستگی به طول توالی در بافرها:} اندازه بافر اتصال میانی تابعی از $P$ است و با افزایش $N$ رشد نمی‌کند.
\end{itemize}

\subsubsection*{۵. ارزیابی تجربی، بنچمارک و توجیه علمی نتایج}
بر اساس آزمون‌های ثبت‌شده در خط لوله ارزیابی پایان‌نامه:
\begin{itemize}
\item \textbf{وضعیت رسمی آزمون کیفی:} \lr{PASS}
\item \textbf{دقت ثبت‌شده روی بنچ‌مارک:} \lr{76.38\%}
\item \textbf{شاخص ارزیابی آماری } \lr{(F1 Score)}\rl{:}  \lr{74.2\%}
\item \textbf{اوج مصرف حافظه گرافیکی } \lr{(Peak VRAM)}\rl{:}  \lr{750.0 MB} (در برابر \lr{۴۲۰۰ MB} مدل مبنا)
\item \textbf{نرخ توان عملیاتی } \lr{(Throughput)}\rl{:}  \lr{37,400} توکن بر ثانیه
\end{itemize}

\textbf{توجیه علمی و علت‌شناسی جذاب نتایج:}
کسب توان بالای \lr{37,400} توکن در ثانیه و محدودسازی حافظه در \lr{750.0 MB} (کاهش \lr{82.1\%}ی نسبت به \lr{۴۲۰۰ MB} پایه) در کنار دقت \lr{76.38\%} نشان می‌دهد که استفاده از گره‌های نهفته واسط، تبادل اطلاعات سراسری را بدون افتادن در دام بن‌بست حافظه تضمین می‌کند.

\subsubsection*{۶. دیاگرام معماری شکیل در لاتک \lr{(TikZ Architecture Flowchart)}}
جهت تبیین دقیق جریان داده تانسوری و نحوه تغییر شکل ماتریس‌ها، نمودار معماری این ماژول با استفاده از بسته قدرتمند \lr{TikZ} در محیط لاتک ترسیم شده است:

\begin{center}
\begin{tikzpicture}[scale=0.82, every node/.style={scale=0.82}, 
    node distance=1.1cm,
    auto,
    block/.style={rectangle, draw=blue!70, fill=blue!12, rounded corners=4pt, text centered, minimum height=0.85cm, font=\small},
    phantom/.style={rectangle, draw=purple!70, fill=purple!15, rounded corners=4pt, text centered, minimum height=0.85cm, font=\small},
    outblock/.style={rectangle, draw=teal!80, fill=teal!15, rounded corners=4pt, text centered, minimum height=0.85cm, font=\small}
]
    \node [block] (input) {تانسور $\mathbf{X}$};
    \node [phantom, right=1.0cm of input] (z) {توکن‌های شبح $\mathbf{Z}$};
    \node [phantom, right=1.0cm of z] (biproj) {پروجکشن دوگانه};
    \node [outblock, right=1.0cm of biproj] (out) {خروجی بازپخش};

    \draw [->, thick, draw=blue!80] (input) -- (z);
    \draw [->, thick, draw=purple!80] (z) -- (biproj);
    \draw [->, thick, draw=purple!80] (biproj) -- (out);
\end{tikzpicture}
\end{center}

\subsubsection*{۷. پیوست تصویری و پرامپت تولید هوش مصنوعی}
برای ترسیم تجسم سه‌بعدی و گرافیک صنعتی این مکانیزم با ابزارهای هوش مصنوعی، پرامپت انگلیسی زیر با دقت تدوین شده است:

\vspace{0.2cm}
\begin{latin}
\footnotesize
\textbf{AI Technical Visualization Prompt:}\\
"A 3D scientific visualization of Phantom attention. Semi-transparent iridescent virtual tokens hovering above a physical data grid, exchanging thin laser threads with dense sequence tokens, futuristic cyberpunk holographic UI, 8k resolution --ar 16:9"
\end{latin}

\begin{figure}[H]
    \centering
    \framebox[0.9\textwidth]{
    \begin{minipage}{0.85\textwidth}
    \vspace{1.8cm}
    \centering
    \textbf{طرح‌واره معماری و جریان تانسوری ماژول \lr{phantom}}\\[0.3cm]\small \lr{(Tensor Flow \& Architectural Schematic for phantom)}\\[0.2cm]
    \small (بر اساس پرامپت فنی فوق تولید و در این کادر جایگزین می‌شود)
    \vspace{1.8cm}
    \end{minipage}
    }
    \caption{شماتیک و جریان داده تانسوری در ماژول \lr{phantom}}
    \label{fig:phantom}
\end{figure}

\vspace{0.5cm}
\hrule
\vspace{0.5cm}
